% pipeline.tex -- the DarwinDiff differentiable parameter-learner pipeline, contrasted
% with the traditional Green's-functions calibration. Scientifically faithful to the
% implementation: three environmental covariates -> a per-cell DINN -> the six bounded
% Carroll-6 parameters -> the differentiable 5-tracer carroll6 box -> an MSE loss vs
% ECCO-Darwin v05, with gradients returned in one backward pass.
% Standalone: `pdflatex pipeline.tex`.  In the paper: \usepackage{standalone}% pipeline.tex -- the DarwinDiff differentiable parameter-learner pipeline, contrasted
% with the traditional Green's-functions calibration. Scientifically faithful to the
% implementation: three environmental covariates -> a per-cell DINN -> the six bounded
% Carroll-6 parameters -> the differentiable 5-tracer carroll6 box -> an MSE loss vs
% ECCO-Darwin v05, with gradients returned in one backward pass.
% Standalone: `pdflatex pipeline.tex`.  In the paper: \usepackage{standalone}% pipeline.tex -- the DarwinDiff differentiable parameter-learner pipeline, contrasted
% with the traditional Green's-functions calibration. Scientifically faithful to the
% implementation: three environmental covariates -> a per-cell DINN -> the six bounded
% Carroll-6 parameters -> the differentiable 5-tracer carroll6 box -> an MSE loss vs
% ECCO-Darwin v05, with gradients returned in one backward pass.
% Standalone: `pdflatex pipeline.tex`.  In the paper: \usepackage{standalone}% pipeline.tex -- the DarwinDiff differentiable parameter-learner pipeline, contrasted
% with the traditional Green's-functions calibration. Scientifically faithful to the
% implementation: three environmental covariates -> a per-cell DINN -> the six bounded
% Carroll-6 parameters -> the differentiable 5-tracer carroll6 box -> an MSE loss vs
% ECCO-Darwin v05, with gradients returned in one backward pass.
% Standalone: `pdflatex pipeline.tex`.  In the paper: \usepackage{standalone}\input{pipeline}.
\documentclass[border=8pt]{standalone}
\usepackage{darwindiff-tikz}
\begin{document}
\begin{tikzpicture}[node distance=6mm and 9mm]

  % ---- DarwinDiff: differentiable autograd pipeline (top) ------------------------
  \node[ddlabel, text=ddIdent, anchor=west] (t1) at (0,1.15) {DarwinDiff --- differentiable (PyTorch autograd)};
  \node[stage, anchor=west]        (env)  at (0,0)      {environment\\[1pt]\ddsub{SST\,$\cdot$\,wind\,$\cdot$\,MLD}};
  \node[stage, right=of env]       (dinn) {DINN\\[1pt]\ddsub{per-cell net}};
  \node[stage, right=of dinn]      (par)  {$\theta$: 6 params\\[1pt]\ddsub{per grid cell}};
  \node[stage, right=of par]       (box)  {differentiable box\\[1pt]\ddsub{carroll6, 5 tracers}};
  \node[stage, right=of box]       (loss) {MSE loss\\[1pt]\ddsub{vs ECCO-Darwin v05}};
  \draw[flow] (env) -- (dinn);
  \draw[flow] (dinn) -- (par);
  \draw[flow] (par) -- (box);
  \draw[flow] (box) -- (loss);
  \draw[grad] (loss.north) to[out=90,in=90]
        node[above=0.5pt, ddtiny, text=ddIdent] {one backward pass $\rightarrow$ all six $\partial\mathcal{L}/\partial\theta$}
        (dinn.north);
  \node[ddtiny, anchor=north] at ([yshift=-1mm]box.south) {5 tracers: DFe, P$_\mathrm{s}$, P$_\mathrm{l}$, POC, PIC};

  % ---- Green's functions: traditional calibration (bottom) ----------------------
  \node[ddlabel, text=ddMuted, anchor=west] (t2) at (0,-2.35) {Green's functions --- traditional (Menemenlis 2005; Carroll 2020)};
  \node[stage, anchor=west]  (g0) at (0,-3.5)  {$\theta$ (6 params)};
  \node[stage, right=of g0]  (g1) {perturb each\\parameter};
  \node[stage, right=of g1, fill=ddAccentBg, draw=ddAccent] (g2) {$N$ full forward\\ECCO-Darwin runs};
  \node[stage, right=of g2]  (g3) {linear\\least-squares};
  \node[stage, right=of g3]  (g4) {global optimum\\[1pt]\ddsub{one value / param}};
  \draw[flow] (g0) -- (g1);
  \draw[flow] (g1) -- (g2);
  \draw[flow] (g2) -- (g3);
  \draw[flow] (g3) -- (g4);
  \node[ddtiny, anchor=north] at ([yshift=-1mm]g2.south) {$\approx$ one full forward run per parameter};

  % ---- the two contrasts --------------------------------------------------------
  \node[verdict=ddIdent, fill=ddIdentBg, anchor=north west, font=\sffamily\scriptsize, text width=5.2cm]
        at ([yshift=-9mm]g0.south west)
        {\textbf{Cost:} 1 backward pass returns all six gradients, vs $\approx\!6$ forward GCM runs per Green's-functions iteration.};
  \node[verdict=ddAccent, fill=ddAccentBg, anchor=north west, font=\sffamily\scriptsize, text width=5.2cm]
        at ([xshift=6mm]$(g0.south west)+(5.2cm,-9mm)$)
        {\textbf{Resolution:} $\theta$ is predicted \emph{per grid cell} from the local environment, vs one global scalar per parameter.};

\end{tikzpicture}
\end{document}
.
\documentclass[border=8pt]{standalone}
\usepackage{darwindiff-tikz}
\begin{document}
\begin{tikzpicture}[node distance=6mm and 9mm]

  % ---- DarwinDiff: differentiable autograd pipeline (top) ------------------------
  \node[ddlabel, text=ddIdent, anchor=west] (t1) at (0,1.15) {DarwinDiff --- differentiable (PyTorch autograd)};
  \node[stage, anchor=west]        (env)  at (0,0)      {environment\\[1pt]\ddsub{SST\,$\cdot$\,wind\,$\cdot$\,MLD}};
  \node[stage, right=of env]       (dinn) {DINN\\[1pt]\ddsub{per-cell net}};
  \node[stage, right=of dinn]      (par)  {$\theta$: 6 params\\[1pt]\ddsub{per grid cell}};
  \node[stage, right=of par]       (box)  {differentiable box\\[1pt]\ddsub{carroll6, 5 tracers}};
  \node[stage, right=of box]       (loss) {MSE loss\\[1pt]\ddsub{vs ECCO-Darwin v05}};
  \draw[flow] (env) -- (dinn);
  \draw[flow] (dinn) -- (par);
  \draw[flow] (par) -- (box);
  \draw[flow] (box) -- (loss);
  \draw[grad] (loss.north) to[out=90,in=90]
        node[above=0.5pt, ddtiny, text=ddIdent] {one backward pass $\rightarrow$ all six $\partial\mathcal{L}/\partial\theta$}
        (dinn.north);
  \node[ddtiny, anchor=north] at ([yshift=-1mm]box.south) {5 tracers: DFe, P$_\mathrm{s}$, P$_\mathrm{l}$, POC, PIC};

  % ---- Green's functions: traditional calibration (bottom) ----------------------
  \node[ddlabel, text=ddMuted, anchor=west] (t2) at (0,-2.35) {Green's functions --- traditional (Menemenlis 2005; Carroll 2020)};
  \node[stage, anchor=west]  (g0) at (0,-3.5)  {$\theta$ (6 params)};
  \node[stage, right=of g0]  (g1) {perturb each\\parameter};
  \node[stage, right=of g1, fill=ddAccentBg, draw=ddAccent] (g2) {$N$ full forward\\ECCO-Darwin runs};
  \node[stage, right=of g2]  (g3) {linear\\least-squares};
  \node[stage, right=of g3]  (g4) {global optimum\\[1pt]\ddsub{one value / param}};
  \draw[flow] (g0) -- (g1);
  \draw[flow] (g1) -- (g2);
  \draw[flow] (g2) -- (g3);
  \draw[flow] (g3) -- (g4);
  \node[ddtiny, anchor=north] at ([yshift=-1mm]g2.south) {$\approx$ one full forward run per parameter};

  % ---- the two contrasts --------------------------------------------------------
  \node[verdict=ddIdent, fill=ddIdentBg, anchor=north west, font=\sffamily\scriptsize, text width=5.2cm]
        at ([yshift=-9mm]g0.south west)
        {\textbf{Cost:} 1 backward pass returns all six gradients, vs $\approx\!6$ forward GCM runs per Green's-functions iteration.};
  \node[verdict=ddAccent, fill=ddAccentBg, anchor=north west, font=\sffamily\scriptsize, text width=5.2cm]
        at ([xshift=6mm]$(g0.south west)+(5.2cm,-9mm)$)
        {\textbf{Resolution:} $\theta$ is predicted \emph{per grid cell} from the local environment, vs one global scalar per parameter.};

\end{tikzpicture}
\end{document}
.
\documentclass[border=8pt]{standalone}
\usepackage{darwindiff-tikz}
\begin{document}
\begin{tikzpicture}[node distance=6mm and 9mm]

  % ---- DarwinDiff: differentiable autograd pipeline (top) ------------------------
  \node[ddlabel, text=ddIdent, anchor=west] (t1) at (0,1.15) {DarwinDiff --- differentiable (PyTorch autograd)};
  \node[stage, anchor=west]        (env)  at (0,0)      {environment\\[1pt]\ddsub{SST\,$\cdot$\,wind\,$\cdot$\,MLD}};
  \node[stage, right=of env]       (dinn) {DINN\\[1pt]\ddsub{per-cell net}};
  \node[stage, right=of dinn]      (par)  {$\theta$: 6 params\\[1pt]\ddsub{per grid cell}};
  \node[stage, right=of par]       (box)  {differentiable box\\[1pt]\ddsub{carroll6, 5 tracers}};
  \node[stage, right=of box]       (loss) {MSE loss\\[1pt]\ddsub{vs ECCO-Darwin v05}};
  \draw[flow] (env) -- (dinn);
  \draw[flow] (dinn) -- (par);
  \draw[flow] (par) -- (box);
  \draw[flow] (box) -- (loss);
  \draw[grad] (loss.north) to[out=90,in=90]
        node[above=0.5pt, ddtiny, text=ddIdent] {one backward pass $\rightarrow$ all six $\partial\mathcal{L}/\partial\theta$}
        (dinn.north);
  \node[ddtiny, anchor=north] at ([yshift=-1mm]box.south) {5 tracers: DFe, P$_\mathrm{s}$, P$_\mathrm{l}$, POC, PIC};

  % ---- Green's functions: traditional calibration (bottom) ----------------------
  \node[ddlabel, text=ddMuted, anchor=west] (t2) at (0,-2.35) {Green's functions --- traditional (Menemenlis 2005; Carroll 2020)};
  \node[stage, anchor=west]  (g0) at (0,-3.5)  {$\theta$ (6 params)};
  \node[stage, right=of g0]  (g1) {perturb each\\parameter};
  \node[stage, right=of g1, fill=ddAccentBg, draw=ddAccent] (g2) {$N$ full forward\\ECCO-Darwin runs};
  \node[stage, right=of g2]  (g3) {linear\\least-squares};
  \node[stage, right=of g3]  (g4) {global optimum\\[1pt]\ddsub{one value / param}};
  \draw[flow] (g0) -- (g1);
  \draw[flow] (g1) -- (g2);
  \draw[flow] (g2) -- (g3);
  \draw[flow] (g3) -- (g4);
  \node[ddtiny, anchor=north] at ([yshift=-1mm]g2.south) {$\approx$ one full forward run per parameter};

  % ---- the two contrasts --------------------------------------------------------
  \node[verdict=ddIdent, fill=ddIdentBg, anchor=north west, font=\sffamily\scriptsize, text width=5.2cm]
        at ([yshift=-9mm]g0.south west)
        {\textbf{Cost:} 1 backward pass returns all six gradients, vs $\approx\!6$ forward GCM runs per Green's-functions iteration.};
  \node[verdict=ddAccent, fill=ddAccentBg, anchor=north west, font=\sffamily\scriptsize, text width=5.2cm]
        at ([xshift=6mm]$(g0.south west)+(5.2cm,-9mm)$)
        {\textbf{Resolution:} $\theta$ is predicted \emph{per grid cell} from the local environment, vs one global scalar per parameter.};

\end{tikzpicture}
\end{document}
.
\documentclass[border=8pt]{standalone}
\usepackage{darwindiff-tikz}
\begin{document}
\begin{tikzpicture}[node distance=6mm and 9mm]

  % ---- DarwinDiff: differentiable autograd pipeline (top) ------------------------
  \node[ddlabel, text=ddIdent, anchor=west] (t1) at (0,1.15) {DarwinDiff --- differentiable (PyTorch autograd)};
  \node[stage, anchor=west]        (env)  at (0,0)      {environment\\[1pt]\ddsub{SST\,$\cdot$\,wind\,$\cdot$\,MLD}};
  \node[stage, right=of env]       (dinn) {DINN\\[1pt]\ddsub{per-cell net}};
  \node[stage, right=of dinn]      (par)  {$\theta$: 6 params\\[1pt]\ddsub{per grid cell}};
  \node[stage, right=of par]       (box)  {differentiable box\\[1pt]\ddsub{carroll6, 5 tracers}};
  \node[stage, right=of box]       (loss) {MSE loss\\[1pt]\ddsub{vs ECCO-Darwin v05}};
  \draw[flow] (env) -- (dinn);
  \draw[flow] (dinn) -- (par);
  \draw[flow] (par) -- (box);
  \draw[flow] (box) -- (loss);
  \draw[grad] (loss.north) to[out=90,in=90]
        node[above=0.5pt, ddtiny, text=ddIdent] {one backward pass $\rightarrow$ all six $\partial\mathcal{L}/\partial\theta$}
        (dinn.north);
  \node[ddtiny, anchor=north] at ([yshift=-1mm]box.south) {5 tracers: DFe, P$_\mathrm{s}$, P$_\mathrm{l}$, POC, PIC};

  % ---- Green's functions: traditional calibration (bottom) ----------------------
  \node[ddlabel, text=ddMuted, anchor=west] (t2) at (0,-2.35) {Green's functions --- traditional (Menemenlis 2005; Carroll 2020)};
  \node[stage, anchor=west]  (g0) at (0,-3.5)  {$\theta$ (6 params)};
  \node[stage, right=of g0]  (g1) {perturb each\\parameter};
  \node[stage, right=of g1, fill=ddAccentBg, draw=ddAccent] (g2) {$N$ full forward\\ECCO-Darwin runs};
  \node[stage, right=of g2]  (g3) {linear\\least-squares};
  \node[stage, right=of g3]  (g4) {global optimum\\[1pt]\ddsub{one value / param}};
  \draw[flow] (g0) -- (g1);
  \draw[flow] (g1) -- (g2);
  \draw[flow] (g2) -- (g3);
  \draw[flow] (g3) -- (g4);
  \node[ddtiny, anchor=north] at ([yshift=-1mm]g2.south) {$\approx$ one full forward run per parameter};

  % ---- the two contrasts --------------------------------------------------------
  \node[verdict=ddIdent, fill=ddIdentBg, anchor=north west, font=\sffamily\scriptsize, text width=5.2cm]
        at ([yshift=-9mm]g0.south west)
        {\textbf{Cost:} 1 backward pass returns all six gradients, vs $\approx\!6$ forward GCM runs per Green's-functions iteration.};
  \node[verdict=ddAccent, fill=ddAccentBg, anchor=north west, font=\sffamily\scriptsize, text width=5.2cm]
        at ([xshift=6mm]$(g0.south west)+(5.2cm,-9mm)$)
        {\textbf{Resolution:} $\theta$ is predicted \emph{per grid cell} from the local environment, vs one global scalar per parameter.};

\end{tikzpicture}
\end{document}
